\scanpage{96}
\mspage{5}
% archive sheet number: 95
on ait $n\gamma=\sum_{0}^{\unsure{n|u}} c_i\beta_i$

On utilise maintenant le théorème de Kummer pour $\mathfrak X/\unsure{\underline D_0|D}$ pour trouver
un \unsure{revêtement} de $\mathfrak X/\underline D_0$ qui ait comme
groupe \deleted{\unsure{d'inertie}} de Galois $J_0$\inserted{$\simeq\mu_n$}, et p.r son
groupe d'inertie en $X_0$ \ldots..

\margin{$\deleted{\unsure{B}}\beta^0=n\gamma+\sum_1^r c_i\beta^{i}$\\
$\unsure{\mathscr L|\mathscr D|\mathcal J}\simeq K^{\otimes n}$}
\diagram{\begin{tikzcd}[ampersand replacement=\&,column sep=small]
\pi_2^t(X_0/D'_0,\bar\xi_0)\arrow[r]\arrow[r,no head,start anchor=south,end anchor=south,to path={-- ++(0,-.28) -| (\tikztotarget)}] \& \mu^\infty_{X_0,t}(\bar\xi_0)\arrow[r] \& \pi_1^t(\mathfrak X/D,\bar\xi_1)\arrow[r] \& \pi_1^t(X_0/D'_0,\bar\xi_0)\arrow[r]\arrow[ll,bend right=28,start anchor=north,end anchor=north east] \& 0
\end{tikzcd}}

$\pi_1^t(\mathfrak X/D,\bar\xi_1)$ comme extension de $\pi_1^t(X_0/D'_0,\bar\xi_0)$
est connu \unsure{entièrement} en termes de
$X_0/D'_0$ et de la classe $\alpha\in H^2(X_0\deleted{\unsure{-D'_0|-D_0}},\mu^\infty_{X_0,\mathrm{tame}})$
[self-intersection de $X_0$ dans $\mathfrak X$] \ldots.

\uline{Exemple I} $X_0$ courbe algébrique \inserted{lisse} sur $k$ \unsure{alg.|sép.} clos.

Alors $\pi_2^t(X_0-D'_0)=0$ . Sauf si $X_0\simeq\mathbb P^1_k$\inserted{, $D'_0=\varnothing$}, en quel
cas $\pi_2^t(X_0,\bar\xi_0)\simeq\deleted{\unsure{A}}\mu^\infty_{X_0,t}(\bar\xi_0)$. L'hom $\mathbb Z\longrightarrow\mu^{\deleted{\unsure{2}}\infty}_t$
dans ce dernier cas correspond \unsure{ex.ct|ex.t} à la
multiplication par $n$ (le degré). \deleted{\unsure{En tout cas}}

1°) \fbox{$X_0-D'_0\ne\mathbb P^1_k$} \deleted{\uline{\unsure{exceptionnel|exception}}} $\pi_1^t(\mathfrak X/D,\bar\xi_{\unsure{1|0}})$ est une extension
\uline{triviale} de $\pi_1^t(X_0-D_0,\bar\xi_0)$ par $\mu^\infty_t(\bar\xi_0)$.
